\documentclass[10pt, aspectratio=169]{beamer}

% --- TEMA Y ESTILO ---
\usetheme{Madrid}
\usecolortheme{whale}

% --- PAQUETES ---
\usepackage[utf8]{inputenc}
\usepackage[spanish, es-noquoting]{babel}
\usepackage{amsmath, amssymb}
\usepackage{siunitx}
\usepackage{booktabs}
\usepackage[american]{circuitikz}
\usetikzlibrary{shapes.geometric, arrows, positioning, babel}

\sisetup{output-decimal-marker = {,}, per-mode = symbol}

% --- INFORMACIÓN DEL CURSO ---
\title[Transitorios y Flyback]{Supresión de Transitorios Inductivos}
\subtitle{Física, Modelado y Protección de Etapas Críticas de Entrada ADC}
\author{M.Sc. Ing. Bernardo Quiroga Turdera}
\institute[UCB Tarija]{Circuitos Electrónicos III (IMT-221) - Sesión 4.2}
\date{Semestre 2-2026}

\begin{document}

\begin{frame}
  \titlepage
\end{frame}

\begin{frame}{Objetivos de la Sesión}
  \begin{itemize}
      \item \textbf{Comprender} la física matemática de los transitorios inductivos (conmutación en cargas reactivas).
      \item \textbf{Modelar y resolver analíticamente} las respuestas temporales del circuito con y sin diodo Flyback.
      \item \textbf{Conectar} la supresión de picos inductivos con la seguridad del módulo analógico del Hito 1.
  \end{itemize}
\end{frame}

\begin{frame}{La Física de la Interrupción Inductiva}
  \begin{block}{El Conflicto Matemático}
      Ecuación Gobernadora de la Bobina:
      \begin{equation*}
          v_L(t) = L \frac{di(t)}{dt}
      \end{equation*}
  \end{block}
  \vspace{0.2cm}
  \textbf{Análisis Conceptual:}
  \begin{itemize}
      \item Un inductor se opone a los cambios bruscos de corriente.
      \item Al pasar un transistor de ON a OFF, la corriente inicial $I_{init}$ es obligada a extinguirse casi instantáneamente ($\Delta t \to 0$).
      \item $\frac{di}{dt}$ tiende a un valor negativo gigantesco $\implies$ Genera un voltaje inducido inverso masivo (\textit{inductive kickback}).
  \end{itemize}
  \alert{\textbf{Consecuencia:}} La tensión se eleva a miles de voltios, destruyendo el semiconductor de potencia por avalancha.
\end{frame}

\begin{frame}{Escenario Crítico SIN Diodo (Circuito Equivalente)}
  \begin{columns}
      \begin{column}{0.4\textwidth}
          \begin{center}
              \begin{circuitikz}[american, scale=0.65, transform shape]
                  \draw (0,0) node[vcc]{$V_{CC}$ (\qty{12}{\volt})} to[L, l=$L$] (0,-2) to[R, l=$R$] (0,-4);
                  \draw (0,-4) to[short, -*] (2,-4) node[right]{$v_{node}$};
                  \draw (0,-4) to[cute open switch, l=Switch] (0,-7) node[ground]{};
                  \draw (2,-4) to[C, l=$C_p$] (2,-7) node[ground]{};
              \end{circuitikz}
          \end{center}
      \end{column}
      \begin{column}{0.6\textwidth}
          El sistema es un \textbf{RLC Altamente Subamortiguado} debido a la pequeñísima capacitancia parásita del transistor ($C_p \approx \qty{100}{\pico\farad}$).
          \begin{equation*}
              L \frac{d^2q(t)}{dt^2} + R_L \frac{dq(t)}{dt} + \frac{q(t)}{C_p} = V_{CC}
          \end{equation*}
          Pico Teórico Máximo (Worst-Case):
          \begin{equation*}
              V_{peak} \approx V_{CC} + I_{init} \sqrt{\frac{L}{C_p}}
          \end{equation*}
      \end{column}
  \end{columns}
\end{frame}

\begin{frame}{Demostración Matemática SIN Diodo}
  \textbf{Ejemplo Práctico:}
  \begin{itemize}
      \item $V_{CC} = \qty{12}{\volt}$, $L = \qty{10}{\milli\henry}$, $R_L = \qty{5}{\ohm}$, $C_p = \qty{100}{\pico\farad}$
      \item Corriente inicial: $I_{init} = \frac{\qty{12}{\volt}}{\qty{5}{\ohm}} = \qty{2.4}{\ampere}$
  \end{itemize}
  \vspace{0.4cm}
  Calculando el pico teórico:
  \begin{equation*}
      V_{peak} \approx \qty{12}{\volt} + \qty{2.4}{\ampere} \times \sqrt{\frac{10^{-2}\text{ H}}{10^{-10}\text{ F}}} = 12 + 2.4 \times 10,000 = \mathbf{24,012\text{ V}}
  \end{equation*}
  \vspace{0.4cm}
  \alert{\textbf{¡24 Kilovoltios en microsegundos!}} Esto destruye inmediatamente el transistor.
\end{frame}

\begin{frame}{Solución de Ingeniería: Diodo Flyback}
  \begin{columns}
      \begin{column}{0.4\textwidth}
          \begin{center}
              \begin{circuitikz}[american, scale=0.6, transform shape]
                  \draw (0,0) node[vcc]{$V_{CC}$ (\qty{12}{\volt})} -- (2,0);
                  \draw (0,0) to[L, l=$L$] (0,-2) to[R, l=$R$] (0,-4);
                  \draw (2,0) to[D, l=$D_1$ Flyback, invert] (2,-4);
                  \draw (0,-4) to[short, -*] (3,-4) node[right]{$v_{node}$};
                  \draw (0,-4) -- (2,-4);
                  \draw (1.5,-4) to[cute open switch, l=Switch] (1.5,-7) node[ground]{};
              \end{circuitikz}
          \end{center}
      \end{column}
      \begin{column}{0.6\textwidth}
          \textbf{Intervalo ON (Transistor Cerrado):}\\
          $v_{node} \approx \qty{0}{\volt}$. Diodo en inversa ($\qty{12}{\volt}$ en cátodo). OFF.
          
          \vspace{0.3cm}
          \textbf{Intervalo Transitorio (Transistor Abre):}\\
          La bobina fuerza la corriente. El voltaje sube hasta polarizar el diodo en directa.
          \begin{equation*}
              v_{node}(t) = V_{CC} + V_{on} \approx \mathbf{\qty{12.7}{\volt}}
          \end{equation*}
          \textit{¡Recorta instantáneamente el pico y salva al transistor!}
      \end{column}
  \end{columns}
\end{frame}

\begin{frame}{Modelo Matemático del Desgaste Exponencial}
  Con el diodo ON, la malla de descarga es un lazo RL cerrado:
  \begin{equation*}
      L \frac{di(t)}{dt} + R_L i(t) + V_{on} = 0
  \end{equation*}
  Solución analítica de la corriente ($\tau = L/R_L$):
  \begin{equation*}
      i(t) = \left( I_{init} + \frac{V_{on}}{R_L} \right) \exp\left( -\frac{t}{\tau} \right) - \frac{V_{on}}{R_L}
  \end{equation*}
  Tiempo exacto de extinción ($i(t_{off})=0$):
  \begin{equation*}
      t_{off} = \tau \ln\left( \frac{V_{CC} + V_{on}}{V_{on}} \right) \implies 2\text{ms} \times 2.899 = \mathbf{\qty{5.798}{\milli\second}}
  \end{equation*}
\end{frame}

\begin{frame}{Conexión Práctica con el Hito 1 del PFI}
  \begin{center}
      \begin{tikzpicture}[node distance=2.8cm, auto, thick, scale=0.7, transform shape]
          \node[draw, rectangle, fill=gray!20, text width=3cm, align=center, minimum height=1.5cm] (motor) {\textbf{Carga Inductiva}\\(Protegida por 1N4007)};
          \node[draw, rectangle, fill=blue!15, text width=3cm, align=center, minimum height=1.5cm, right of=motor, node distance=3.5cm] (shunt) {\textbf{Resistencia Shunt}\\(Sujeción Negativa 1N4148)};
          \node[draw, rectangle, fill=orange!15, text width=3cm, align=center, minimum height=1.5cm, right of=shunt, node distance=3.5cm] (zener) {\textbf{Zener Limitador}\\(Sujeción Positiva 5.1V)};
          \node[draw, rectangle, fill=green!15, text width=2.5cm, align=center, minimum height=1.5cm, right of=zener, node distance=3.2cm] (adc) {\textbf{Entrada ADC}\\(Rango Seguro)};
          \draw[->, >=stealth, line width=1.5pt] (motor) -- (shunt);
          \draw[->, >=stealth, line width=1.5pt] (shunt) -- (zener);
          \draw[->, >=stealth, line width=1.5pt] (zener) -- (adc);
      \end{tikzpicture}
  \end{center}
  
  \textbf{Protección Complementaria del ADC:}\\
  Incluso con el Flyback apagando el pico masivo de \qty{24}{\kilo\volt} en el motor, las inductancias de los cables inyectan ruido rápido a la placa lógica.
  \begin{itemize}
      \item \textbf{Zener 5.1V:} Sujeta sobrepicos por arriba de \qty{5.1}{\volt}.
      \item \textbf{1N4148 ($t_{rr} \le \qty{4}{\nano\second}$):} Limitador ultra rápido. Drena picos negativos al instante protegiendo el silicio del ADC.
  \end{itemize}
\end{frame}

\begin{frame}{Taller Computacional y Dinámica ABET}
  \begin{center}
      Ejecutaremos el script de Python para visualizar el análisis transitorio y validar el modelo matemático (Pilar 1).
  \end{center}
  
  \vspace{0.4cm}
  \textbf{Auditoría y Pensamiento Crítico:}\\
  Al cambiar la inductancia $L$ de $\qty{10}{\milli\henry}$ a $\qty{50}{\milli\henry}$ en vivo, el script debe predecir un tiempo de descarga (Rueda Libre) sustancialmente más largo. Esto lo \textbf{deberán comprobar visualmente con el osciloscopio} en el laboratorio del viernes.
\end{frame}

\end{document}
